\documentclass[12pt]{article}

\usepackage{amsmath}
\usepackage{geometry}
\geometry{a4paper, margin=2.5cm}

\begin{document}

\title{Brukerveiledning -- FOPDT Interaktiv Modell i Google Colab}
\author{}
\date{}
\maketitle

\section*{Hvordan åpne notatboken}
\begin{enumerate}
    \item Klikk på lenken læreren har gitt deg for å åpne notatboken i Google Colab.
    \item Du trenger ikke Google-konto for å kjøre den, men det anbefales.
\end{enumerate}

\section*{Slik bruker du modellen}
\begin{enumerate}
    \item Last opp målefilen \texttt{pid\_test.txt} dersom notatboken ber om det.
    \item Juster sliderne:
    \begin{itemize}
        \item $A$ – amplituden på steget
        \item $T$ – tidskonstant
        \item $L$ – død tid
        \item $y_0$ – startnivå
    \end{itemize}
    \item Observer hvordan den røde modellkurven oppdateres.
    \item Tilpass modellen slik at den passer best mulig til de blå måledataene.
    \item Noter de endelige verdiene for $A$, $T$, $L$ og $y_0$ som står under grafen.
\end{enumerate}

\section*{Beregn PID-parametre}
Bruk verdiene du fant til å beregne PID-parametrene.  
Dette skal du gjøre for hånd.

\[
K = \frac{A}{\Delta u}
\]

\[
K_p = \frac{T}{K \cdot (\lambda + L)}
\]

\[
T_i = T
\]

\[
T_d = \frac{L \cdot T}{L + \lambda}
\]

\section*{Lagre figur eller fil}
\begin{itemize}
    \item Klikk \textbf{Last ned PNG} for å lagre grafen.
    \item Du kan lagre din egen kopi av notatboken via:
    \begin{itemize}
        \item \texttt{File → Download .ipynb}, eller
        \item \texttt{File → Save a copy in Drive}
    \end{itemize}
\end{itemize}

\section*{Viktig}
\begin{itemize}
    \item Du kan ikke lagre endringer tilbake til lærerens GitHub.
    \item Du kan jobbe fritt i din egen kopi uten å ødelegge originalen.
\end{itemize}

\end{document}